% Source: https://artofproblemsolving.com/community/c3h2454557_geometric_proof_of_cauchyschwartz_inequality
% Stable identity: aops-topic-2454557-q01
Let $a,b,x_1,y_1$ be real numbers with $(a,b)\ne(0,0)$.
\begin{enumerate}
\item[(a)] The distance from the point $P(x_1,y_1)$ to the line $\ell:\ ax+by=0$ is $\dfrac{|ax_1+by_1|}{\sqrt{a^2+b^2}}$. Use this and the fact that $\ell$ passes through the origin $O$ to prove the Cauchy-Schwarz inequality
\[(ax_1+by_1)^2\le\left(a^2+b^2\right)\left(x_1^2+y_1^2\right),\]
and find when equality holds.
\item[(b)] A student writes the distance as $\dfrac{ax_1+by_1}{\sqrt{a^2+b^2}}$, concludes $ax_1+by_1\le\sqrt{x_1^2+y_1^2}\sqrt{a^2+b^2}$, and then squares, saying ``both sides are nonnegative''. What is wrong with this step, and how is it repaired?
\end{enumerate}
